\subsection{Signed-source transfer and sharper complete-source cell caps}
\label{sec:signed-source}
The local route can compare the signed source directly, without separating
its multiplier into independent magnitudes. For two holomorphic fields
with common $q=q(t,x)>0$, $r=\log q$, let $d_j$ bound their jet differences
for $j=0,1,2$ and let common $P_1,P_2$ bound their jets and straight
interpolation. On a common amplitude interval $[m,M]$ put
\begin{align}
 \mathcal S_F&=A_s(|F|)\{r_t|F_z|^2+r_x\Re(F_{zz}\overline{F_z})\},\notag\\
 A_s(v)&=\frac{q(v^2-q)}{\pi(q+v^2)^3},\quad
 A_s'(v)=\frac{4qv(2q-v^2)}{\pi(q+v^2)^4},\notag\\
 A_s''(v)&=\frac{4q(2q^2-17qv^2+5v^4)}{\pi(q+v^2)^5}.
 \label{eq:signed-radial}
\end{align}
The factorization follows from \eqref{eq:adaptiveSource}, and ordinary
differentiation gives the derivatives. Let $a_s=\sup_{[m,M]}|A_s|$,
$a_{s,1}=\sup_{[m,M]}|A_s'|$. The first maximum is attained among the
endpoints and $\sqrt{2q}$; the second among the endpoints and
$\sqrt{q(17\pm\sqrt{249})/10}$, retaining only points inside the
interval. These are the stationary points from the displayed derivatives;
interior zeros of the absolute-valued functions are minima, so endpoints
also cover degenerate intervals.

\begin{proposition}[Direct signed-source comparison]\label{prop:signed-source}
With the common frozen caps above,
\begin{align}
 |\mathcal S_H-\mathcal S_F|\le E_{\mathcal S}
 ={}&a_s\{2|r_t|P_1d_1+|r_x|(P_1d_2+P_2d_1)\}\notag\\
 &+a_{s,1}(|r_t|P_1^2+|r_x|P_1P_2)d_0.
 \label{eq:signed-source-bound}
\end{align}
No heat equation is required. For every fixed spatial derivative used,
\begin{equation}
 |\partial^\alpha K*(\mathcal S_H-\mathcal S_F)|
                    \le|K_\alpha|*E_{\mathcal S}.
 \label{eq:signed-source-filter}
\end{equation}
\end{proposition}
\begin{proof}
Write the bracket in \eqref{eq:signed-radial} as $T_F$.
Product comparison bounds its difference by
$2|r_t|P_1d_1+|r_x|(P_1d_2+P_2d_1)$ and its modulus by
$|r_t|P_1^2+|r_x|P_1P_2$. The real gradient norm of the radial map
$A_s(|F|)$ is $|A_s'|$ along the field segment, including its smooth
limit at zero. Telescope $A_sT$ to prove the first inequality. Transfer
the derivative of the entire product to the fixed kernel for the second.
Thus the varying $r_t,r_x$ stay inside the full-support integral; no
regulator derivative is omitted or replaced by a center value.
\end{proof}
The row is a positive linear map in $d_0,d_1,d_2$. Freeze its coefficients
with the total errors before separately mapping interpolation, modes,
complete tail and nodal arithmetic. Its validity extends to fields with
heat residuals because the signed source was defined independently.
Note~36's optional direct $W_t,C,C_t$ envelopes instead assume exact heat
fields; a joined or saddle field needs the residual terms in
\eqref{eq:cj-direct-time} and Note~36 AS11. They are not imported without
that premise \cite{n36}.

For source curvature the complete signed derivative polynomial is
\begin{equation}
 \Phi_n''=\pi^{-1/4}n^{-1/2}e^{-A_n}A_n^{5/4}P_2(A_n),
 \qquad P_2(A)=-75+330A-224A^2+32A^3.
 \label{eq:signed-source-curvature}
\end{equation}
On $[l,l+d]$ its cubic Bernstein coefficients are
\[
 P_2(l),\quad P_2(l)+dP_2'(l)/3,\quad
 P_2(l)+2dP_2'(l)/3+d^2P_2''(l)/6,\quad P_2(l+d).
\]
Their largest absolute value bounds $|P_2|$ continuously. Multiply by
$\pi^{-1/4}n^{-1/2}m_{5/4}(l,l+d)$ and sum retained modes.
For a cap sum through $N_B$, $n_0=N_B+1$ and
$A_0=\pi n_0^2e^{4a}\ge17/4$, the full omitted-mode cap is
\begin{equation}
 \frac{\pi^{-1/4}n_0^{-1/2}e^{-A_0}
 (75A_0^{5/4}+330A_0^{9/4}+224A_0^{13/4}+32A_0^{17/4})}
 {1-(1+1/n_0)^8e^{-\pi(2n_0+1)e^{4a}}},
 \label{eq:signed-curvature-tail}
\end{equation}
provided the ratio is below one. Monomial decrease and the maximum
mode-power exponent eight prove this tail as before.
A second valid cell bound, for midpoint $c$ and width $h$, is
$|\Phi''(c)|+(h/2)\sup|\Phi'''|$, with complete source and derivative
tails retained. The recurrence gives
$P_3=-375+3270A-4232A^2+1440A^3-128A^4$; positive monomial caps
give the full third-derivative tail, once $A_0\ge21/4$,
\[
 \frac{\pi^{-1/4}n_0^{-1/2}e^{-A_0}
 \sum_{j=0}^4 |p_{3,j}|A_0^{5/4+j}}
 {1-(1+1/n_0)^{10}e^{-\pi(2n_0+1)e^{4a}}},
\]
where $p_{3,j}$ are the displayed coefficients of $P_3$ and the ratio
is required to be below one. Taking the minimum with the earlier complete-source
cap remains valid. The finite-source mode count is separate from this
cap evaluation count.

For nonnegative time ceiling and $Y$, the weights
$u^je^{tu^2+Yu}$ are convex on $u\ge0$. Integrate the nonnegative
linear-interpolation remainder kernel and hats against their secant
chord majorants. The resulting degree-at-most-three polynomial integrals
bound the weights continuously; for constant $W$ they give $h^3W/12$
and $hW/2$. The cubic counterparts are \eqref{eq:hermite-chords}.
This pays a chord-majorant inequality, not floating endpoint or summation
arithmetic. Other time domains need their own weight bound.
The adaptive mesh heuristic and numerical comparisons remain in Note~36;
no new numerical study is part of this integration.
